\section{Introduction}
\label{sec:introduction}

This is the second paper in a series developing a deterministic brane-lattice substrate hypothesis.
Paper~I established the substrate model --- a 4D cubic lattice embedded with codimension zero in a
4D Euclidean ambient, governed by a hyperelastic action built from the induced metric --- and derived
the linearized wave-branch structure with characteristic speeds
$c_T^2=(1-\alpha)k_s a^2/m$ and $c_L^2=k_s a^2/m$.
The present paper develops two consistency bridges from that foundation.

\paragraph{What this paper imports.}
Both bridges import from Paper~I the following interface equations:
the near-identity metric expansion (\Cref{eq:metric-near-identity}),
the linearized branch speeds (\Cref{eq:wave-speeds-derived}),
the long-wavelength dispersion (\Cref{eq:linear-dispersion}),
and the quasi-static in-brane equilibrium equations (\Cref{subsec:gravity-channel}).
No new substrate assumptions are introduced here.

\paragraph{Lorentz bridge.}
The cubic lattice is manifestly anisotropic at the lab-frame level: the longitudinal wave speed
along $[100]$ differs from the speed along $[111]$ by roughly 7.5\% at $\alpha=0.2$.
The Lorentz bridge (\Cref{sec:lorentz}) argues that this anisotropy cancels for inside observers
whose rods, clocks, and signal cones are made of the same substrate material.
The cancellation is not tuned; it is a structural consequence of the substrate-universality
principle: the same stiffness tensor that sets branch speed also determines the elastic response
of any rod-shaped excitation.
Residual birefringence --- the splitting of effective signal speeds that survives at finite
wavelength and amplitude --- is the primary falsifiable prediction.

\paragraph{Gravity bridge.}
The induced metric $g_{ij} = h_\star^2\delta_{ij} + \partial_i u\,\partial_j u$ contains a
quadratic source term in the transverse gradient.
When the in-brane displacement field $\vec\xi$ is free to respond, quasi-static balance forces
a slowly varying contraction pattern inward toward any region of localized transverse excitation.
The gravity bridge (\Cref{sec:gravity}) identifies this as the substrate origin of a long-range
attractive channel.
The bridge is geometrically fixed; the open quantitative task is to derive a Newtonian-limit
Poisson equation from the coarse-grained contraction field and express the effective coupling
$G_{\text{eff}}$ in substrate parameters.

\paragraph{Why these two bridges are paired.}
Both bridges operate through the same geometric mechanism --- the Pythagorean lengthening of
links by transverse displacement --- and in the fully symmetric 4D treatment they are two faces
of one rule: spacelike excitation lengthens timelike links (time dilation) and timelike
excitation lengthens spacelike links (gravity channel / length contraction).
Pairing them in one paper makes this unification explicit and avoids treating the two phenomena
as independent postulates.

\paragraph{What this paper does not claim.}
The gravity bridge is established only at the level of a geometrically fixed contraction channel.
This paper does \emph{not} claim a completed quantitative reduction of that channel to a Newtonian
limit, nor a value of the effective coupling $G$ derived from substrate parameters; that reduction
is posed as the open quantitative task of \Cref{sec:gravity}.

\paragraph{Paper organization.}
\Cref{sec:interfaces} lists the interface equations imported from Paper~I.
\Cref{sec:lorentz} develops the Lorentz bridge: wave cone, dual-observer argument, analogue-metric
viewpoint, birefringence as falsifiable signature, and open problems.
\Cref{sec:gravity} develops the gravity bridge: contraction-channel derivation, Newtonian-limit
open problem, and simulation test specification.